%Conversational system is becoming an indispensable component in a variety of AI applications and acts as an interactive interface provided to users to improve user experience. Language understanding is crucial for conversational systems to provide appropriate responses to users. Intent classification is usually the first step of language understanding in these systems. Its primary goal is to identify diverse intentions behind user utterances and it is usually formalized as a classification problem. However, intent classes defined during training are inevitably inadequate to cover all of the user intents at the inference stage due to flexibility and randomness of the user utterances. Hence, out-of-scope (or unknown) intent classification is introduced to address this problem and it aims at building a classifier that can accurately identify supported (or known) intent classes while detecting out-of-scope (or unknown) intents that are unrelated to supported intents.


Conversational system is becoming an indispensable component in a variety of AI applications and acts as an interactive interface provided to users to improve user experience. Language understanding is essential for conversational systems to provide appropriate responses to users, and intent detection is usually the first step of language understanding. The primary goal is to identify diverse intentions behind user utterances, which is often formalized as a classification task. However, intent classes defined during training are inevitably inadequate to cover all possible user intents at the test stage due to the diversity and randomness of user utterances. Hence, out-of-scope (or unknown) intent detection is essential, which aims to develop a model that can accurately identify known (seen in training) intent classes while detecting the out-of-scope classes that are not encountered during training.

Due to the practical importance of out-of-scope intent detection, recent efforts have attempted to solve this problem by developing effective intent classification models. In general, previous works approach this problem by learning \emph{decision boundaries} for known intents and then using some confidence measure to distinguish known and unknown intents. For examples, LMCL~\citep{DBLP:conf/acl/LinX19} learns the decision boundaries with a margin-based optimization objective, and SEG~\citep{yan-etal-2020-unknown} assumes the known intent classes follow the distribution of mixture of Gaussians. After learning the decision boundaries, an off-the-shell outlier detection algorithm such as LOF~\citep{10.1145/335191.335388} is commonly employed to derive confidence scores  \cite{yan-etal-2020-unknown,DBLP:conf/emnlp/ShuXL17, DBLP:conf/acl/LinX19, DBLP:conf/iclr/HendrycksG17}.
If the confidence score of a test sample is lower than a predefined threshold, it is identified as an outlier. 

However, it may be problematic to learn decision boundaries solely based on the training examples of known intent classes. First, if there are sufficient training examples, the learned decision boundaries can be expected to generalize well on known intent classes, but not on the unknown. Therefore, extra steps are required in previous methods, such as using an additional outlier detection algorithm at the test stage or adjusting the confidence threshold by cross-validation. On the other hand, if there are not sufficient training examples, the learned boundaries may not generalize well on both known and unknown intents. As a result, these methods often underperform when not enough training data is given. Hence, it is important to provide learning signals of unknown intents at the training stage to overcome these limitations.

%is the key point of removing the constraints of previous methods. 

%In this paper, our goal is to remove these limitations and propose a simple yet effective method for out-of-scope intent detection.  We take a brave step to build pseudo out-of-scope examples directly to aid the training process. leading to the flexibility of selceting them.

In contrast to previous works, we adopt a different approach by explicitly modeling the distribution of unknown intents. Particularly, we construct a set of pseudo out-of-scope examples to aid the training process. We hypothesize that in the semantic feature space, real-world outliers can be well represented in two types: ``\emph{hard}'' outliers that are geometrically close to the inliers and ``\emph{easy}'' outliers that are distant from the inliners. For the ``hard'' ones, we construct them in a self-supervised manner by forming convex combination of the features of inliers from different classes. 
For the ``easy'' ones, the assumption is that they are very unrelated to the known intent classes, so they can be used to simulate the randomness and diversity of user utterances. They can be easily constructed using public datasets. For example, in our experiments, we randomly collect sentences from datasets of other NLP tasks such as question answering and sentiment analysis as open-domain outliers. 

In effect, by constructing pseudo outliers for the unknown class during training, we form a consistent $(K+1)$ classification task ($K$ known classes + 1 unknown class) for both training and test. Our model can be trained with a cross-entropy loss and directly applied to test data for intent classification and outlier detection without requiring any further steps.
As shown in Figure~\ref{fig:tsne} (better view in color and enlarged), our method can learn better utterance representations, which make each known intent class more compact and push the outliers away from the inliers. Our main contributions are summarized as follows.
\begin{itemize}
%\item We propose a novel model that is simple, efficient, and effective for out-of-scope intent detection by constructing pseudo outliers. We form a consistent (K+1) classification task for the training and test process to bridge the gap between fitting to training data and generalizing to test data.

\item We propose a novel out-of-scope intent detection approach by matching training and test tasks to bridge the gap between fitting to training data and generalizing to test data.

\item We propose to efficiently construct two types of pseudo outliers by using a simple self-supervised method and leveraging publicly available auxiliary datasets.

\item We conduct extensive experiments on four real-world dialogue datasets to demonstrate the effectiveness of our method and perform a detailed ablation study.
\end{itemize}